% Job posting: https://ca.linkedin.com/jobs/view/embedded-software-designer-%E2%80%93-processing-algorithms-at-kepler-communications-inc-4454747734
\documentclass[11pt]{article}
\usepackage[left=0.7in,top=0.7in,right=0.7in,bottom=0.7in]{geometry}
\usepackage{enumitem}
\usepackage[hidelinks]{hyperref}
\usepackage{titlesec}
\usepackage{array}
\usepackage{setspace}
\usepackage{needspace}
\usepackage[T1]{fontenc}
\usepackage{newtxtext,newtxmath}
\ifdefined\pdfgentounicode
  \input{glyphtounicode}
  \pdfgentounicode=1
\fi
\setstretch{1.04}
\pagenumbering{gobble}
\titleformat{\section}{\Large\scshape}{}{4em}{}[\vspace{-0.7em}\rule{\linewidth}{0.4pt}\vspace{-0.4em}]
\titlespacing*{\section}{0pt}{1em}{0.4em}
\newenvironment{cvrole}[5]{%
  \par\noindent\textbf{\large #1} | \textit{\small #2, #3}\hfill \textit{\small #4 - #5}\par
  \begin{itemize}[leftmargin=2em, itemsep=-0.3em, topsep=-0.5em]%
}{%
  \end{itemize}\par\noindent\mbox{}\par\vspace{0.1em}%
}
\newcommand{\cvName}{Nishal Stanislaus Silva}
\newcommand{\cvEmail}{nishal.silva@hotmail.com}
\newcommand{\cvPhone}{+1 613 200 2030}
\newcommand{\cvCity}{Toronto, ON}
\newcommand{\cvSite}{https://nishal.xyz}
\newcommand{\cvLinkedIn}{https://www.linkedin.com/in/nishal-silva}
\newcommand{\cvGitHub}{https://github.com/ns2max}
\newcommand{\cvStatus}{Canadian Permanent Resident, Eligible to work in Canada without sponsorship}
\newcommand{\cvTagline}{Embedded Systems Engineer | Real-Time DSP, Signal Processing \& Constrained Hardware}
\begin{document}
\pagestyle{empty}
\begin{center}
    {\LARGE \scshape \textbf{\cvName}}{\large \textit{, PhD}}\\[1pt]
    {\small \cvTagline}\\[1pt]
    {\footnotesize\textit{\cvStatus}}\\[1pt]
    {\small
    \href{mailto:\cvEmail}{\cvEmail} \,|\, \cvPhone \,|\, \cvCity
    \,|\, \href{\cvSite}{nishal.xyz} \,|\, \href{\cvLinkedIn}{linkedin.com/in/nishal-silva}
    \,|\, \href{\cvGitHub}{github.com/ns2max}}
\end{center}

\section*{Professional Summary}
Embedded signal-processing engineer with a PhD in real-time algorithm design under hard latency, compute and power budgets. MSc thesis and Asilomar 2018 paper apply the S-transform to time-frequency extraction from noisy continuous signals. As a postdoctoral researcher, designed and shipped a real-time pattern-detection algorithm running at 14 ms inference on Raspberry Pi 4 (F1 0.76, 31.4\% CPU), 74x faster than the DTW baseline it replaced -- algorithm design under genuinely constrained embedded hardware, published in JAES 2026. Built MIRaaS, a UDP edge-to-server offload architecture with characterized latency (IEEE IS2 2026), directly relevant to networked real-time signal transmission. At McGill, profiled compute and power draw on a self-powered, battery-constrained embedded platform to validate real-time feasibility -- power-budget engineering that transfers directly to spacecraft-class hardware constraints. Seeks to bring embedded C++17 algorithm design and constrained-hardware discipline to small-satellite communications systems.

\section*{Core Competencies}
\textbf{Embedded Signal Processing \& Algorithms:} STFT, MFCC/LFCC/PLP/GFCC, S-transform (Stockwell transform), CQT, chroma features, onset/beat tracking, YIN pitch estimation, time-frequency analysis of noisy real-world signals, real-time pattern-detection algorithm design, benchmark and ablation design for accuracy/latency/CPU trade-offs. \\
\textbf{Constrained Hardware Engineering:} Real-time C++17 systems under hard latency budgets, ARM CPU and power profiling, embedded Linux, IMU/I2C sensor integration, PLC interfacing, ADC/DAC hardware integration, self-powered/battery-constrained platform validation, networked real-time signal transmission (UDP offload architectures).

\section*{Technical Stack}
C++17 (expert, real-time engines), Python, Raspberry Pi, ARM, embedded Linux, I2C, PLC integration, HiFiBerry ADC/DAC, AWS, Docker, CI/CD, Git, pytest.

\section*{Education}
\noindent\textbf{PhD, Information Engineering and Computer Science} -- University of Trento, Italy \hfill 2021 - 2025 \\
Defended Jan 2025. Thesis: Embedded Real-time Musical Pattern Detection for Smart Musical Instruments. \\[4pt]
\noindent\textbf{MSc, Telecommunication and Electronic Engineering} -- Sheffield Hallam University, UK \hfill 2016 - 2018 \\
Thesis: On Musical Onset Detection via the S-Transform (published, IEEE Asilomar 2018) -- time-frequency signal-processing method for noisy continuous signals. \\[4pt]
\noindent\textbf{BEng (Hons), Electronic Engineering} -- Sheffield Hallam University, UK \hfill 2013 - 2014

\section*{Research and Work Experience}

\begin{cvrole}{Independent Researcher}{Self-directed}{Toronto, Canada}{Jan 2026}{Present}
\item Two papers accepted for publication (Smart Drums, JAES; MIRaaS, IEEE IS2 2026).
\item Maintain Nebula, a C++17 audio-feature-extraction library with per-feature ARM/Raspberry Pi latency benchmarks (13 stars).
\item Continued embedded real-time DSP and algorithm-benchmarking work using AI-assisted development workflows including Claude Code.
\end{cvrole}

\begin{cvrole}{Postdoctoral Researcher}{University of Trento}{Trento, Italy}{Jan 2025}{Dec 2025}
\item Owned the full pipeline for streaming audio + IMU/gesture data on the EU Horizon EIC Pathfinder project MUSMET: acquisition, DSP/feature engineering, training/eval, edge deployment on Raspberry Pi 4, monitoring.
\item Designed and shipped a real-time pattern-detection algorithm: 14 ms inference on Raspberry Pi 4, F1 0.76, 31.4\% CPU, 74x faster than the 1038 ms DTW baseline it replaced (JAES 2026).
\item Designed frozen-protocol benchmark suites and ablations quantifying accuracy/latency/CPU trade-offs for embedded deployment decisions.
\item Built MIRaaS, a UDP edge-to-server offload architecture with characterized latency for networked real-time signal transmission (IEEE IS2 2026, accepted).
\end{cvrole}

\begin{cvrole}{Visiting Researcher}{McGill University (IDMIL / CIRMMT)}{Montreal, Canada}{May 2024}{Aug 2024}
\item Built a self-powered smart electric guitar: BNO055 IMU (I2C) + HiFiBerry (audio ADC/DAC hardware) + Raspberry Pi 4.
\item Fused IMU and audio signals via a hierarchical FSM, F1 0.78 on a 500-event corpus (IEEE IS2 2025).
\item Profiled compute and power draw to validate real-time feasibility on self-powered, battery-constrained embedded hardware.
\end{cvrole}

\begin{cvrole}{Visiting Researcher}{University of Visual and Performing Arts}{Colombo, Sri Lanka}{Jan 2024}{Mar 2024}
\item Conducted a human-centred user study of musicians' perception of smart instruments with embedded pattern detection (IEEE I3DA 2025).
\end{cvrole}

\begin{cvrole}{Technical Consultant}{Forestpin (Pvt) Ltd}{Colombo, Sri Lanka}{Aug 2020}{Feb 2021}
\item Built ML and statistical pipelines for financial forensics, compliance monitoring and risk detection on large structured enterprise data.
\end{cvrole}

\begin{cvrole}{Research Engineer}{MAS Holdings (Pvt) Ltd}{Colombo, Sri Lanka}{Jan 2015}{Jul 2020}
\item Delivered end-to-end computer-vision and IoT systems for manufacturing at scale (5.5 years), lifting operational efficiency up to +300\%.
\item Ran real-time IoT ETL across production sites with C++/Python inference deployed on embedded hardware; introduced CI/CD and code review; mentored engineers.
\item Built automated multilingual care-label QC (OpenCV, conveyor + PLC reject integration), cutting inspection time by 99.5\%.
\end{cvrole}

\section*{Publications}
Silva, N., Weeraddana, P.C. \& Fischione, C. (2018). On Musical Onset Detection via the S-Transform. IEEE Asilomar Conference on Signals, Systems, and Computers -- lead author, time-frequency signal-extraction method for noisy continuous signals. Silva, N. \& Turchet, L. (2026). Real-Time Audio Pattern Detection for Smart Musical Instruments. J. Audio Eng. Soc. 74(3). Silva, N. \& Turchet, L. (2026, in press). Smart Drums. J. Audio Eng. Soc. Silva, N. \& Turchet, L. (2026, accepted). MIRaaS. IEEE International Symposium on the Internet of Sounds (IS2).

\section*{Highlights}
MIDI Innovation Awards 2023 Finalist ("Hot Licks", Prototypes \& Non-Commercial Software category, covered by Sound On Sound and MusicRadar). GMOA recognition (Government Medical Officers' Association, Sri Lanka, 2020) for a COVID-19 remote vitals monitoring deployment built and shipped for hospital ward use. Track record spans lead-author signal-processing research, real-time embedded algorithm design validated on constrained hardware, and production hardware integration (PLC, IMU, ADC/DAC) at manufacturing scale.

\end{document}
